\slajdid{10-s041}
\begin{frame}[label=10-s041]
\gusttekst\setlength{\abovedisplayskip}{3pt}\setlength{\belowdisplayskip}{3pt}\setlength{\abovedisplayshortskip}{2pt}\setlength{\belowdisplayshortskip}{2pt}\begin{columns}[c,onlytextwidth]\column{.29\textwidth}\centering\begin{tikzpicture}[x=.62cm,y=.8cm,font=\sitneoznake,>=Latex]\ctikzset{bipoles/length=.6cm}\draw(-1.8,1)node[ocirc]{}node[left]{$V_{I1}$}--(-1,1);\begin{scope}[shift={(0,1)}]\draw(-1,0)--(-.3,0);\draw(-.3,-.3)--(-.3,.3)--(.3,0)--cycle;\draw(.3,0)--(1,0);\draw(.3,-.3)--(.3,.3);\end{scope}\draw(1,1)--(1.4,1);\node[below]at(0,0.85){D1};\draw[<->](-.5,1.55)--(.5,1.55);\draw(-.5,1.4)--(-.5,1.7);\draw(.5,1.4)--(.5,1.7);\node[above]at(0,1.7){$+\ V_{D1}$};\draw[->](-1.7,1.35)--(-1.1,1.35)node[midway,above]{$I_{D1}$};\draw(-1.8,-1)node[ocirc]{}node[left]{$V_{I2}$}--(-1,-1);\begin{scope}[shift={(0,-1)}]\draw(-1,0)--(-.3,0);\draw(-.3,-.3)--(-.3,.3)--(.3,0)--cycle;\draw(.3,0)--(1,0);\draw(.3,-.3)--(.3,.3);\end{scope}\draw(1,-1)--(1.4,-1);\node[below]at(0,-1.15){D2};\draw[<->](-.5,-0.44999999999999996)--(.5,-0.44999999999999996);\draw(-.5,-0.6)--(-.5,-0.30000000000000004);\draw(.5,-0.6)--(.5,-0.30000000000000004);\node[above]at(0,-0.30000000000000004){$+\ V_{D2}$};\draw[->](-1.7,-0.65)--(-1.1,-0.65)node[midway,above]{$I_{D2}$};\draw(1.4,1)--(1.4,-1);\draw(1.4,0)--(2.8,0)node[ocirc]{}node[right]{$V_O$};\draw(2.1,0)to[R,l=$R$](2.1,-1.8)node[ground]{};\draw[->](1.6,-.5)--(1.6,-1.1);\node[left]at(1.6,-1.55){$I_R$};\end{tikzpicture}
\column{.69\textwidth}
Zbog simetrije istu situaciju bi dobili i da je $V_{I2}\gg V_{\gamma D}$ dok je $V_{I1}$ i dalje manje od $V_{\gamma D}$. Tada bi vodila dioda D2, dioda D1 bi bila zakočena, inverzno polarizovana, a izraz za izlazni napon bi bio $V_O=V_{I2}-V_D$.
\par\medskip Na potpuno identičan način u slučaju $V_{I1}=V_{I2}\gg V_{\gamma D}$ se vidi da će obe diode i D1 i D2 voditi. Dele struju $I_R$.
\par\medskip Znači na ulazu $(V_L,V_L)$ na izlazu $V_L$. Ako je na bilo kojem ulazu $V_H$ biće visok napon i na izlazu. Logička funkcija je dvoulazna ILI funkcija. Uočiti da će se kolo isto ponašati i ako dodamo na isti način D3, D4, …, odnosno da na lak način možemo povećati broj ulaza u kolo.
\end{columns}\par\medskip\centering\begin{tikzpicture}[x=.8cm,y=.6cm,font=\sitneoznake,>=Latex]\draw[->](-.2,0)--(3.2,0)node[below]{$V_I$};\draw[->](0,-.2)--(0,2.6)node[left]{$V_O$};\draw(.4,0)--(2.6,2.2);\node[below]at(.4,0){$V_D$};\end{tikzpicture}
\par\medskip\raggedright
\textcolor{DEcrvena}{Ovo kolo ne može biti logičko kolo pošto mu karakterisitka prenosa ne odgovara, ali obezbeđuje jednostavn višelaznu logičku ILI funkciju}
\end{frame}
